\providecommand{\Ind}{\operatorname{Ind}}

\chapter{George Lusztig (2022)}
\index{Lusztig, George}

\begin{quote}
\it ``George Lusztig has been the dominant figure in the representation theory of reductive groups over finite fields for the past fifty years. His development of character sheaves, the theory of cells in Coxeter groups, and the geometric construction of quantum groups have transformed the field and created deep connections between representation theory, algebraic geometry, and combinatorics.'' --- Wolf Prize Citation, 2022
\end{quote}

\section{Early Life and Career}

George Lusztig was born on May 20, 1946, in Timi\c{s}oara, Romania. He studied at the University of Bucharest, receiving his Ph.D.\ in 1968 under the supervision of Alexandru Froda. He emigrated to the United States in 1969 and has been a professor at the Massachusetts Institute of Technology since 1978. He was awarded the Wolf Prize in Mathematics in 2022.

\section{Main Contributions}

George Lusztig (born 1946) is a Romanian--American mathematician whose work has fundamentally shaped representation theory, algebraic geometry, and combinatorics. He was awarded the Wolf Prize in Mathematics in 2022.

The Wolf Prize citation reads: ``for his groundbreaking contributions to representation theory and to related areas.''

Lusztig's core contributions include:
\begin{itemize}
    \item \textbf{Character Sheaves:} The development of the theory of character sheaves, a geometric approach to constructing characters of finite groups of Lie type, providing a uniform description of all irreducible characters.
    \item \textbf{Lusztig--Kazhdan Theory\index{Kazhdan-Lusztig theory}:} The classification of irreducible representations of affine Hecke algebras and the establishment of the Kazhdan--Lusztig conjectures, later proved by Beilinson--Bernstein and Brylinski--Kashiwara.
    \item \textbf{Cells in Coxeter Groups:} The theory of left, right, and two-sided cells in Coxeter groups, with applications to Hecke algebras, representation theory, and total positivity.
    \item \textbf{Canonical Bases\index{canonical basis}:} The construction of canonical bases for quantum groups and their classical limits, providing a deep connection between quantum groups and combinatorics.
    \item \textbf{Representations of Finite Groups of Lie Type:} The classification and character formula for all irreducible representations of reductive groups over finite fields.
    \item \textbf{Geometric Quantum Groups:} The Beilinson--Lusztig--MacPherson geometric realization of quantum $\GL_n$ using perverse sheaves on quiver varieties.
    \item \textbf{Total Positivity:} The theory of totally positive matrices and its generalization to arbitrary Lie groups, with applications to cluster algebras.
\end{itemize}

\section{Origin of the Problem}

\subsection{Representations of Finite Groups of Lie Type}

Finite groups of Lie type (such as $\GL(n, \F_q)$, $\SL(n, \F_q)$, etc.) constitute most of the finite simple groups. Understanding their irreducible representations and characters was a central problem in group theory. The Deligne--Lusztig theory (1976) provided a construction of representations using l-adic cohomology of Deligne--Lusztig varieties, but the classification was not complete.

Lusztig's character sheaves provided a uniform geometric framework that completed the classification.

\subsection{Hecke Algebras and Kazhdan--Lusztig Conjectures}

The Hecke algebra\index{Hecke algebra} $H_q(W)$ of a Coxeter group $W$ is a deformation of the group algebra. The Kazhdan--Lusztig polynomials $P_{x,w}(q)$ encode subtle combinatorial information about the Bruhat order. The Kazhdan--Lusztig conjectures (proved later) relate these polynomials to the characters of Verma modules in representation theory.

\subsection{Canonical Bases}

Quantum groups $U_q(\mathfrak{g})$ have a distinguished basis (the canonical basis) with remarkable positivity properties. Lusztig's construction of this basis using perverse sheaves established deep connections between quantum groups, algebraic geometry, and combinatorics.

\section{How It Was Solved and the Intuitions/Tools Needed}

\subsection{Character Sheaves}

A character sheaf on a reductive group $G$ over a finite field is a complex of sheaves on $G$ that is an eigenobject for the Fourier transform. Lusztig showed that the characteristic functions of character sheaves (evaluated at Frobenius fixed points) give the irreducible characters of $G(\F_q)$. The approach uses:
\begin{enumerate}
    \item The geometry of the group $G$ and its flag variety.
    \item The theory of perverse sheaves and l-adic cohomology.
    \item The Fourier--Deligne transform on the group.
\end{enumerate}

\subsection{Kazhdan--Lusztig Theory}

The Kazhdan--Lusztig polynomials $P_{x,w}(q)$ are defined recursively using the Hecke algebra relations. The key insight: these polynomials encode the local intersection cohomology of Schubert varieties in the flag variety. This connection between combinatorics (Hecke algebras) and geometry (Schubert varieties) was revolutionary.

\subsection{Canonical Bases}

Lusztig constructed the canonical basis for $U_q^+(\mathfrak{g})$ using the geometry of quiver varieties and the theory of perverse sheaves. The basis elements correspond to irreducible components of certain Lagrangian varieties, and their structure constants have remarkable positivity properties.

\section{Mathematical Theory}

\subsection{Character Sheaves}

\begin{definition}[Character Sheaf]
Let $G$ be a connected reductive group over $\F_q$. A \textbf{character sheaf} on $G$ is a simple perverse sheaf $\mathcal{A}$ on $G$ (up to shift) that is a Fourier--Deligne eigenobject: there exists a local system on the set of semisimple conjugacy classes such that $\mathcal{A}$ is an eigenobject under the Fourier transform.
\end{definition}

\begin{theorem}[Lusztig, 1980s]
The irreducible characters of the finite group $G(\F_q)$ are given by the characteristic functions of character sheaves evaluated at Frobenius fixed points. More precisely, for each character sheaf $\mathcal{A}$, the function:
\[
\chi_{\mathcal{A}}(g) = \sum_i (-1)^i \tr(\Frob, \mathcal{H}^i(\mathcal{A})_g)
\]
is (up to sign and scaling) an irreducible character of $G(\F_q)$.
\end{theorem}

\subsection{Kazhdan--Lusztig Polynomials}

\begin{definition}[Hecke Algebra]
The \textbf{Hecke algebra} $H_q(W)$ of a Coxeter group $(W,S)$ has basis $\{T_w: w \in W\}$ with relations:
\[
T_s T_w = \begin{cases}
T_{sw} & \text{if } \ell(sw) > \ell(w) \\
q T_{sw} + (q-1) T_w & \text{if } \ell(sw) < \ell(w)
\end{cases}
\]
\end{definition}

\begin{definition}[Kazhdan--Lusztig Polynomials]
For $x, w \in W$ with $x \leq w$ in the Bruhat order, the \textbf{Kazhdan--Lusztig polynomial} $P_{x,w}(q)$ is defined by:
\[
C_w = \sum_{x \leq w} (-1)^{\ell(w) - \ell(x)} q^{\ell(w)/2 - \ell(x)} P_{x,w}(q^{-1}) T_x
\]
where $C_w$ are the Kazhdan--Lusztig basis elements.
\end{definition}

\begin{theorem}[Kazhdan--Lusztig, 1979; proved 1981]
The Kazhdan--Lusztig polynomial $P_{x,w}(q)$ evaluated at $q = 1$ gives the multiplicity of the Verma module $M_x$ in the Jordan--H\"older series of the Verma module $M_w$. Equivalently, $P_{x,w}(1)$ is the multiplicity of the simple module $L_x$ in $M_w$.
\end{theorem}

\subsection{Canonical Bases}

\begin{definition}[Canonical Basis]
For a semisimple Lie algebra $\mathfrak{g}$, the \textbf{canonical basis} $\mathbf{B}$ of $U_q^+(\mathfrak{g})$ is a basis with the properties:
\begin{enumerate}
    \item $\mathbf{B}$ is invariant under the bar involution.
    \item $\mathbf{B}$ is compatible with the PBW decomposition.
    \item Structure constants are in $\N[q, q^{-1}]$ (positivity).
\end{enumerate}
\end{definition}

\begin{theorem}[Lusztig, 1990]
The canonical basis $\mathbf{B}$ is constructed geometrically using perverse sheaves on quiver varieties. Each basis element corresponds to an irreducible component of a Lagrangian variety in a quiver variety.
\end{theorem}

\subsection{Cells in Coxeter Groups}

\begin{definition}[Cells]
The \textbf{left cells} of a Coxeter group $W$ are the equivalence classes under the relation: $x \sim_L y$ if there exist Kazhdan--Lusztig basis elements $C_w$ such that $C_x$ and $C_y$ appear in the same left ideal. Similarly define right cells and two-sided cells.
\end{definition}

\begin{theorem}[Lusztig, 1980s]
The two-sided cells of a Weyl group $W$ are in bijection with the unipotent classes of the corresponding algebraic group. Moreover, each left cell gives rise to a representation of $W$ (the cell representation).
\end{theorem}

\subsection{Construction of Canonical Bases via Kashiwara Operators}

\begin{definition}[Kashiwara Operators]
For a quantum group $U_q(\mathfrak{g})$ with Chevalley generators $E_i, F_i, K_i$, the \textbf{Kashiwara operators} $\tilde{E}_i$ and $\tilde{F}_i$ are defined on the crystal basis $\mathcal{B}$ of a highest weight module $V(\lambda)$ by:
\[
\tilde{E}_i(b) = \begin{cases} E_i \cdot b / [n+1]_q & \text{if } E_i^{(n+1)} \cdot b \neq 0 \text{ and } E_i^{(n)} \cdot b = 0 \text{ with } n = \varepsilon_i(b) \\ 0 & \text{otherwise} \end{cases}
\]
\[
\tilde{F}_i(b) = \begin{cases} F_i \cdot b / [n+1]_q & \text{if } F_i^{(n+1)} \cdot b \neq 0 \text{ and } F_i^{(n)} \cdot b = 0 \text{ with } n = \varphi_i(b) \\ 0 & \text{otherwise} \end{cases}
\]
where $[n]_q = (q^n - q^{-n})/(q - q^{-1})$ is the quantum integer, and $E_i^{(n)} = E_i^n/[n]_q!$.
\end{definition}

\begin{definition}[Crystal Graph]
The \textbf{crystal graph} of $V(\lambda)$ is the directed graph $(\mathcal{B}, E)$ where the vertices are elements of $\mathcal{B}$ and there is an arrow $b \to b'$ labeled $i$ if $\tilde{F}_i(b) = b'$. The crystal graph encodes the combinatorial structure of the representation.
\end{definition}

\begin{theorem}[Kashiwara, 1991; Lusztig, 1990]
The canonical basis $\mathbf{B}$ of $U_q^+(\mathfrak{g})$ can be constructed using Kashiwara operators and crystal graphs. The construction proceeds as follows:
\begin{enumerate}
    \item For each dominant weight $\lambda$, the crystal $\mathcal{B}(\lambda)$ has a unique highest weight element $u_\lambda$.
    \item The crystal basis $\mathcal{B}$ of $U_q^-(\mathfrak{g})$ is the disjoint union $\bigsqcup_\lambda \mathcal{B}(\lambda)$.
    \item Each element $b \in \mathcal{B}$ corresponds to a canonical basis element $\mathbf{b} \in \mathbf{B}$ via the crystal-to-basis map.
    \item The bar involution $\overline{\cdot}$ on $U_q^+(\mathfrak{g})$ acts on the crystal by $\overline{b} = b$ for all $b \in \mathcal{B}$ (up to normalization).
\end{enumerate}
The resulting basis $\mathbf{B} = \{\mathbf{b}: b \in \mathcal{B}\}$ is the canonical basis.
\end{theorem}

\begin{remark}
The crystal graph provides a purely combinatorial description of the canonical basis. The structure constants of the canonical basis (the coefficients in the product $\mathbf{b}_1 \cdot \mathbf{b}_2 = \sum c_{b_1, b_2}^{b_3} \mathbf{b}_3$) lie in $\N[q, q^{-1}]$, which is Lusztig's celebrated positivity property. This was proved using the geometric construction via perverse sheaves on quiver varieties.
\end{remark}

\subsection{Kazhdan--Lusztig Polynomials and Intersection Cohomology}

\begin{definition}[Recursive Definition of KL Polynomials]
The Kazhdan--Lusztig polynomials $P_{x,w}(q)$ for $x, w \in W$ with $x \leq w$ are defined recursively by:
\begin{enumerate}
    \item $P_{x,x}(q) = 1$ for all $x \in W$.
    \item If $s \in S$ and $\ell(sw) < \ell(w)$ (so $sw < w$), then:
    \[
    P_{x,w}(q) = \begin{cases}
    P_{sx, sw}(q) & \text{if } \ell(sx) > \ell(x) \\
    q \cdot P_{sx, sw}(q) + (q^{1/2} - q^{-1/2}) \cdot P_{x, sw}(q) - q \cdot P_{sx, w}(q) & \text{if } \ell(sx) < \ell(x)
    \end{cases}
    \]
    \item The polynomial $P_{x,w}(q)$ has degree at most $\frac{1}{2}(\ell(w) - \ell(x) - 1)$.
    \item When $\ell(w) - \ell(x)$ is odd, $P_{x,w}(q)$ is divisible by $q^{1/2}$.
\end{enumerate}
\end{definition}

\begin{theorem}[Kazhdan--Lusztig, 1979]
The Kazhdan--Lusztig polynomial $P_{x,w}(q)$ satisfies $P_{x,w}(1) = \dim L_x^{(w)}$, where $L_x^{(w)}$ is the simple subquotient of the Verma module $M_w$ of highest weight $-w(\rho) - \rho$. Equivalently:
\[
[M_w: L_x] = P_{x,w}(1)
\]
for all $x \leq w$ in the Bruhat order.
\end{theorem}

\begin{theorem}[Intersection Cohomology Connection; Kazhdan--Lusztig, 1979]
Let $X_w$ be the Schubert variety in the flag variety $G/B$ corresponding to $w \in W$. Then for $x \leq w$:
\[
P_{x,w}(q) = \sum_{i=0}^{d} \dim \mathcal{H}^i_{c}(X_w \cap X_w^{\circ}, \Q_l) \cdot q^{i/2}
\]
where $\mathcal{H}^i_c$ denotes the intersection cohomology groups and $X_w^{\circ}$ is the Bruhat cell corresponding to $w$. In particular:
\[
\dim \mathcal{H}^{2i}(X_w) = \sum_{x \leq w} [M_w: L_x] = \sum_{x \leq w} P_{x,w}(1)
\]
\end{theorem}

\begin{remark}
The intersection cohomology of Schubert varieties is Poincar\'e duality, and the Kazhdan--Lusztig polynomials encode the Betti numbers of these varieties. This geometric interpretation was the key to proving the Kazhdan--Lusztig conjectures via the Beilinson--Bernstein--Deligne--Gabber localization theorem and the Brylinski--Kashiwara proof using D-modules.
\end{remark}

\subsection{Total Positivity and the Unipotent Variety}

\begin{definition}[Total Positivity]
A matrix $A \in \GL_n(\R)$ is \textbf{totally positive} if all its minors are positive. More generally, $A$ is \textbf{totally nonnegative} if all minors are nonnegative. The set of totally positive matrices forms a semigroup under multiplication.
\end{definition}

\begin{theorem}[Lusztig, 1992; Total Positivity in Reductive Groups]
Let $G$ be a split reductive group over $\R$ with Weyl group $W = (s_1, \ldots, s_n)$. For each reduced expression $w_0 = s_{i_1} \cdots s_{i_N}$ of the longest element $w_0 \in W$, Lusztig constructs:
\begin{enumerate}
    \item Unipotent elements $u_1, \ldots, u_N \in G(\R)$ (the \textbf{Lusztig parameterization}).
    \item A map $\mathfrak{L}_{w_0}: \R^N \to G(\R)$ given by $\mathfrak{L}_{w_0}(t_1, \ldots, t_N) = u_1(t_1) \cdots u_N(t_N)$.
    \item The image $\mathfrak{L}_{w_0}(\R_{>0}^N)$ is the set of \textbf{totally positive} elements of $G$.
\end{enumerate}
\end{theorem}

\begin{theorem}[Lusztig's Canonical Basis for the Unipotent Variety]
The totally nonnegative part of the unipotent variety $U$ of a reductive group $G$ admits a cell decomposition:
\[
U_{\geq 0} = \bigsqcup_{w \in W} C_w
\]
where $C_w$ are cells parameterized by elements of the Weyl group, and each cell is homeomorphic to $\R^{\ell(w)}$. The closure relations are encoded by the Kazhdan--Lusztig polynomials:
\[
\overline{C_w} = \bigsqcup_{x \leq w} C_x
\]
\end{theorem}

\begin{remark}
Lusztig's total positivity program has deep connections to cluster algebras (Fomin--Zelevinsky). The totally nonnegative part of the unipotent variety is related to the space of phylogenetic trees and to Teichm\"uller theory. The parameterization of totally positive matrices using reduced expressions of $w_0$ provides a combinatorial description of the cell structure.
\end{remark}

\begin{remark}
The unipotent variety $U$ of a reductive group $G$ is the set of elements $g \in G$ such that $g$ is conjugate to a unipotent matrix (all eigenvalues equal to $1$). Lusztig's work shows that the intersection of the unipotent variety with the totally nonnegative part of $G$ has a rich combinatorial structure governed by cells parameterized by the Weyl group.
\end{remark}

\section{Impact on Science and Technology}

\subsection{In Mathematics}

\begin{enumerate}
    \item \textbf{Representation Theory:} Character sheaves provide the most complete and uniform approach to representations of finite groups of Lie type.
    \item \textbf{Algebraic Geometry:} The connections between perverse sheaves and representation theory are fundamental.
    \item \textbf{Combinatorics:} Kazhdan--Lusztig polynomials, cells, and canonical bases are central objects.
    \item \textbf{Quantum Groups:} Lusztig's canonical basis is the definitive construction.
    \item \textbf{Total Positivity:} Applications to cluster algebras and Teichm\"uller theory.
\end{enumerate}

\subsection{In Physics}

\begin{enumerate}
    \item \textbf{Quantum Field Theory:} Canonical bases and quantum groups appear in integrable systems and conformal field theory.
    \item \textbf{Statistical Mechanics:} The Yang--Baxter equation and quantum groups.
\end{enumerate}

\section{Frequently Asked Questions}

\subsection*{Q1: What are character sheaves?}
A geometric approach to the representation theory of finite groups of Lie type. Character sheaves are perverse sheaves on a reductive group whose characteristic functions give the irreducible characters of the finite group.

\subsection*{Q2: What are Kazhdan--Lusztig polynomials?}
Polynomials $P_{x,w}(q)$ associated to pairs of elements in a Coxeter group, defined using the Hecke algebra. They encode multiplicities in Verma modules and the intersection cohomology of Schubert varieties.

\subsection*{Q3: What is the canonical basis?}
A distinguished basis for quantum groups with positivity properties, constructed geometrically by Lusztig. It is fundamental to the theory of quantum groups and has connections to combinatorics and total positivity.

\subsection*{Q4: What are cells in Coxeter groups?}
Equivalence classes in a Coxeter group defined via the Kazhdan--Lusztig basis. They give rise to representations of the Weyl group and correspond to unipotent classes.

\subsection*{Q5: What should an undergraduate study to understand Lusztig's work?}
Representation theory (Lie algebras, finite groups), algebraic geometry (sheaves, cohomology), and combinatorics (Coxeter groups, Hecke algebras). Recommended: \textit{Representation Theory} (Humphreys), \textit{Introduction to Quantum Groups} (Kassel), and \textit{Lectures on Algebraic Geometry} (Harris).

\section{Related Chapters}
See also Chapter~28 (Tits) on buildings and groups of Lie type and Chapter~26 (Thompson) on finite simple groups.

\section*{Exercises for the Reader}
\addcontentsline{toc}{section}{Exercises}

\begin{enumerate}
    \item [I] Compute the Kazhdan--Lusztig polynomials for the symmetric group $S_3$.
    \item [B] Show that the Hecke algebra $H_q(S_n)$ specializes to $\C[S_n]$ when $q = 1$.
    \item [I] Compute the left cells of $S_3$ and describe the corresponding cell representations.
    \item [B] Show that the canonical basis for $U_q(\mathfrak{sl}_2)$ consists of elements $E^{(n)} = E^n/[n]_q!$.
    \item [A] Prove that a character sheaf on $GL(2)$ gives the characters of $GL(2, \F_q)$.
    \item [B] Show that the Kazhdan--Lusztig polynomials for the dihedral group $D_4$ are all $1$.
    \item [A] Compute the irreducible characters of $\SL(2, \F_q)$ using Deligne--Lusztig theory.
\end{enumerate}

\section*{Timeline}
\addcontentsline{toc}{section}{Timeline}

\begin{tabularx}{\linewidth}{|p{1.5cm}|>{\raggedright\arraybackslash}X|}
\hline
\textbf{Year} & \textbf{Event} \\ \hline
1946 & Born in Timi\c{s}oara, Romania \\ \hline
1964 & Moves to Switzerland, then the United States \\ \hline
1971 & Ph.D.\ from Princeton (advisor: William Browder) \\ \hline
1974 | Professor at MIT \\ \hline
1979 | Kazhdan--Lusztig polynomials \\ \hline
1980s | Character sheaves \\ \hline
1985 | Cole Prize of the AMS \\ \hline
1990 | Canonical bases for quantum groups \\ \hline
1998 | Brouwer Medal \\ \hline
2002 | Steele Prize of the AMS \\ \hline
2022 | Wolf Prize in Mathematics \\ \hline
\end{tabularx}

\section*{Glossary}
\addcontentsline{toc}{section}{Glossary}

\begin{description}
    \item[Canonical basis] A distinguished basis of $U_q^+(\mathfrak{g})$ with positivity properties.
    \item[Cell] An equivalence class in a Coxeter group defined via the Kazhdan--Lusztig basis.
    \item[Character sheaf] A perverse sheaf on a reductive group whose characteristic function gives an irreducible character.
    \item[Coxeter group] A group with presentation generated by involutions with specified braid relations.
    \item[Hecke algebra] A deformation of the group algebra of a Coxeter group.
    \item[Kazhdan--Lusztig polynomial] A polynomial encoding intersection cohomology of Schubert varieties.
    \item[Perverse sheaf] An object in the derived category of constructible sheaves satisfying support conditions.
    \item[Quiver variety] A moduli space of representations of a quiver.
    \item[Schubert variety] A subvariety of the flag variety defined by Bruhat order conditions.
\end{description}

\section{Citations}\subsection*{Primary Sources}
\begin{enumerate}
    \item D.\ Kazhdan and G.\ Lusztig, \textit{Representations of Coxeter groups and Hecke algebras}, Invent.\ Math.\ 53 (1979), 165--184.
    \item D.\ Kazhdan and G.\ Lusztig, \textit{Schubert varieties and Poincar\'e duality}, Proc.\ Symp.\ Pure Math.\ 36 (1980), 185--203.
    \item G.\ Lusztig, \textit{Characters of Reductive Groups over a Finite Field}, Ann.\ Math.\ Studies 107, Princeton Univ.\ Press, 1984.
    \item G.\ Lusztig, \textit{Introduction to Quantum Groups}, Birkh\"auser, 1993.
    \item G.\ Lusztig, \textit{Canonical bases arising from quantized enveloping algebras}, J.\ Amer.\ Math.\ Soc.\ 3 (1990), 447--498.
\end{enumerate}

\subsection*{Expository Works}
\begin{enumerate}
    \item R.\ W.\ Carter, \textit{Finite Groups of Lie Type}, Wiley, 1985.
    \item S.\ Riche and G.\ Williamson, \textit{Tilting modules and the p-canonical basis}, Ast\'erisque 397, 2018.
    \item A.\ A.\ Beilinson and J.\ Bernstein, \textit{A proof of the Kazhdan--Lusztig conjecture}, \textit{Advances in Soviet Math.} 16 (1993), 1--50.
\end{enumerate}

